\chapter{Mathematical content}

\section{Number systems} \label{sec:nums}

The file \texttt{number-systems.scm} declares five number systems as
class constants and installs their arithmetic axioms in the
library.  All declarations take effect at load time.

\subsection{The standard number domains}

\begin{itemize}
\item \texttt{nn} ($\mathbb{N}$) Non-negative integers $\{0, 1, 2, \ldots\}$
\item \texttt{zz} ($\mathbb{Z}$) Integers
\item \texttt{qq} ($\mathbb{Q}$) Rationals
\item \texttt{rr} ($\mathbb{R}$) Reals
\item \texttt{cc} ($\mathbb{C}$) Complex numbers
\end{itemize}

Sethood axioms (\texttt{nn-is-set} through \texttt{cc-is-set}) and the inclusion chain
\[
  \mathbb{N} \subseteq \mathbb{Z} \subseteq \mathbb{Q} \subseteq \mathbb{R} \subseteq \mathbb{C}
\]
\begin{sloppypar}are provided as axioms \texttt{nn-subset-zz}, \texttt{zz-subset-qq},
\texttt{qq-subset-rr}, \texttt{rr-subset-cc}.\end{sloppypar}

\subsection{Arithmetic operations}

Arithmetic operations are shared across all number systems: there is one
\texttt{+}, one \texttt{*}, one \texttt{-}, and so on.  Membership
determines which system an argument belongs to; closure axioms guarantee
the result stays in the appropriate system.  For example,
\texttt{forall([x in nn, y in nn], x + y in nn)} is the
closure axiom for \texttt{+} on \texttt{nn}.

\begin{center}
\begin{tabular}{lll}
\toprule
System & Operations & Distinguished constants \\
\midrule
\texttt{nn} &
  \texttt{+}, \texttt{*}, \texttt{succ}, \texttt{<=} &
  \texttt{0}, \texttt{1} \\[3pt]
\texttt{zz} &
  \texttt{+}, \texttt{*}, \texttt{-} (unary), \texttt{<=} &
  \texttt{0}, \texttt{1} \\[3pt]
\texttt{qq} &
  \texttt{+}, \texttt{*}, \texttt{-}, \texttt{recip}, \texttt{<=} &
  \texttt{0}, \texttt{1} \\[3pt]
\texttt{rr} &
  \texttt{+}, \texttt{*}, \texttt{-}, \texttt{recip}, \texttt{abs}, \texttt{<=} &
  \texttt{0}, \texttt{1} \\[3pt]
\texttt{cc} &
  \texttt{+}, \texttt{*}, \texttt{-}, \texttt{recip}, \texttt{conjugate} &
  \texttt{0}, \texttt{1} \\
\bottomrule
\end{tabular}
\end{center}

\begin{remark}\label{remark:27135-34942-842417}
%
The infix notation \texttt{x \^{} n} is the preferred way to write
exponentiation; it parses to the same internal form as
\texttt{power(x, n)}.  The typing axiom is the natural closure form,
matching the two-argument application syntax:
\[
  \forall x\,\forall n.\;
    (x \in \mathtt{cc} \land n \in \mathtt{nn})
    \;\Rightarrow\; \texttt{power}(x,n) \in \mathtt{cc}
  \qquad\text{(axiom \texttt{power-typing-nonneg})}
\]
The defining axioms are \texttt{power-zero} ($x^0 = 1$) and
\texttt{power-succ} ($x^{\mathrm{succ}(n)} = x \cdot x^n$ for
$n \in \mathtt{nn}$).  A conditional equation \texttt{power-neg}
covers negative integer exponents: $x^{-n} = \mathrm{recip}(x^n)$
for $n \in \mathtt{nn}$, $x \neq 0$; but this does not extend the
\texttt{fun} typing to $\mathtt{zz}$, since $0^{-1}$ is undefined.
Complex exponentiation $x^z$ for $z \in \mathtt{cc}$ will be defined
later via the power series for \texttt{exp}.
%
\end{remark}
%
\begin{remark}\label{remark:27135-34995-146430}
%
In particular, $0^0 = 1$.  This is the mathematically correct value
for combinatorial and power-series purposes, as Knuth argues, and it is
what MIT~Scheme's exact arithmetic returns.  GNU/Maxima raises an error
here.  IMPS treated $0^0$ as undefined, which is logically sound but
forces theorems such as the binomial theorem to carry an explicit
$x \neq 0$ or $y \neq 0$ hypothesis to exclude the special case;
this system avoids that overhead by defining $0^0 = 1$.
%
\end{remark}

The arithmetic decision procedure \texttt{vnb-do-arith} (and the
\texttt{(arith)} command) evaluates \texttt{x \^{} n} on ground
numeric arguments.

\noindent
Integer literals are constants directly: \texttt{0}, \texttt{1},
\texttt{7}, \texttt{-3}, etc.\ are elements of \texttt{zz} (and of
\texttt{qq}, \texttt{rr}, \texttt{cc} by inclusion); non-negative integer
literals are additionally elements of \texttt{nn}.  No subscripted
forms (\texttt{0\_ZZ}, \texttt{1\_RR}, etc.) are used.  In MIT~Scheme
numeric literals are first-class objects, so this representation is
natural.

Integer literals must be in canonical decimal form: no leading zeros
(other than \texttt{0} itself), so \texttt{007} and \texttt{000075} are
rejected by the parser.  \texttt{-0} normalises to \texttt{0}.

The ordinal operators \texttt{ord-le}, \texttt{ord-lt}, and
\texttt{succ\_ord} are named apart from the numeric \texttt{<=} and
\texttt{succ}.  The first two were written \texttt{<=\_ORD} and
\texttt{<\_ORD} until 2026-08-24; those spellings could not be tokenized
(\texttt{read-op} consumes the \texttt{<} and stops, and \texttt{\_}
cannot begin a token), so no formula mentioning them could be retyped
from its own printed form.  \texttt{succ\_ord} keeps its subscript: it
begins with a letter and reads back.

\subsection{Axioms installed}

Each number system comes with: a sethood axiom, inclusion from the
previous system, closure of each operation, and algebraic identities
appropriate to its structure.

\begin{itemize}
\item \texttt{nn}: closed under \texttt{+} and \texttt{*}; constants
  \texttt{0} (additive unit) and \texttt{1} (multiplicative unit);
  commutativity, associativity, and distributivity for \texttt{+} and
  \texttt{*}.  Unary \texttt{-} and \texttt{recip} are not total here.
\item \texttt{zz}: $\mathtt{nn}$ axioms plus closure under unary
  \texttt{-} (additive inverse); $x + (-x) = 0$.
\item \texttt{qq}: $\mathtt{zz}$ axioms plus \texttt{recip} on nonzero
  elements; ordering \texttt{<=} compatible with \texttt{+} and \texttt{*}.
\item \texttt{rr}: $\mathtt{qq}$ axioms plus \texttt{abs} satisfying the
  triangle inequality, with \texttt{abs($x$) = 0} iff \texttt{$x$ = 0}.
\item \texttt{cc}: $\mathtt{rr}$ axioms plus \texttt{conjugate}, an
  involution compatible with $+$ and $\cdot$, with
  $a \cdot \overline{a} \in \mathbb{R}$ and $a \cdot \overline{a} \geq 0$.
  No ordering on $\mathtt{cc}$.
\end{itemize}

\noindent The abstract structural names---commutative monoid,
commutative ring, ordered field---and the corresponding officialised
ring instances are deferred to Section~\ref{sec:basic-rings}.  At this
point we are simply listing the operations and identities each
domain carries.

Because there is a single \texttt{+}, \texttt{*}, etc.\ shared across all
systems, no compatibility axioms between adjacent systems are needed:
there is nothing to be compatible with.

\begin{remark}
Each of \texttt{nn}, \texttt{zz}, and \texttt{qq} carries an infinite
family of membership axioms, one per numeral appearing in the library.
When a positive-integer numeral $n$ is introduced, it is implicitly
axiomatized as a iterated sum $1 + 1 + \cdots + 1$ ($n$ times), so
$n \in \mathtt{nn}$.  For \texttt{zz}, negative numerals are similarly
axiomatized as the additive inverse of such sums, giving $-n \in
\mathtt{zz}$.  These axioms are generated on demand: entering a numeral
literal in a formula automatically extends the library with the
corresponding membership fact.
\end{remark}

\subsection{Example use in a proof}

\begin{VNBcode}
; Prove: 0 is in RR
(sp (make-wff-from-string "0 in RR"))
(ta 'rr-zero-in)                ; introduce theorem (0 in rr)
(ass)                           ; goal is now in assumptions

; Prove: for reals a, b, (a + b) is in RR
(sp (make-wff-from-string "a in RR and b in RR implies a + b in RR"))
(di)                              ; assume (a in rr) and (b in rr)
(ai "a in RR and b in RR")       ; split into (a in rr) and (b in rr)
(ta 'rr-add-closed)               ; closure axiom: forall([a, b], ...)
(inst "forall([a, b], a in rr and b in rr implies a + b in rr)" 'a)
(inst "forall([b], a in rr and b in rr implies a + b in rr)" 'b)
(bc "a in rr and b in rr implies a + b in rr")
(ass)
; => Proof complete.
\end{VNBcode}

%% =========================================================
\section{Ordinals}
\label{sec:ordinals}

\subsection{The class \texttt{ord}}

\texttt{ord} is a primitive class constant, declared in the same style as
\texttt{nn}, \texttt{rr}, and \texttt{cc}: no explicit construction of
ordinals is needed.  The class is characterised entirely by its axioms.

\begin{itemize}
\item \texttt{ord} the class of all ordinals (a proper class)
\item \texttt{$\alpha$ in ord} $\alpha$ is an ordinal
\end{itemize}

\subsection{The ordering predicate}

Ordinals are equipped with a primitive ordering predicate \texttt{ord-le}
(and its strict companion \texttt{ord-lt}) that is independent of set
membership.  The ordering has no required connection to $\in$.

\begin{itemize}
\item \texttt{ord-le($\alpha$, $\beta$)} \quad $\alpha \leq \beta$ in the ordinal ordering
\item \texttt{ord-lt($\alpha$, $\beta$)} \quad $\alpha < \beta$ (strict)
\item \texttt{succ\_ord($\alpha$)} \quad successor ordinal of $\alpha$
\end{itemize}

\subsection{Axioms}

The following axioms are installed in \texttt{ordinals.scm}.

\begin{itemize}
\item \texttt{burali-forti}
  $\mathtt{ord} \notin \mathtt{set}$: \texttt{ord} is a proper class.
  Every \emph{set} of ordinals has a supremum in \texttt{ord}
  (axiom \texttt{sup-ord-in}); the supremum need not lie in the set itself
  (e.g.\ $\sup \mathtt{nn} = \omega \notin \mathtt{nn}$).
  If \texttt{ord} were a set, its supremum $\alpha$ would be an ordinal,
  hence $\mathit{succ\_ord}(\alpha) \in \mathtt{ord}$ would strictly
  exceed $\alpha$---contradicting that $\alpha$ bounds \texttt{ord}.
\item \texttt{ord-le-closure}
  $\alpha \leq_O \beta \;\Rightarrow\; \alpha \in \mathtt{ord} \land \beta \in \mathtt{ord}$
\item \texttt{ord-le-refl}
  $\alpha \in \mathtt{ord} \;\Rightarrow\; \alpha \leq_O \alpha$
\item \texttt{ord-le-antisymm}
  $\alpha \leq_O \beta \land \beta \leq_O \alpha \;\Rightarrow\; \alpha = \beta$
\item \texttt{ord-le-trans}
  $\alpha \leq_O \beta \land \beta \leq_O \gamma \;\Rightarrow\; \alpha \leq_O \gamma$
\item \texttt{ord-le-total}
  $\alpha,\beta \in \mathtt{ord} \;\Rightarrow\; \alpha \leq_O \beta \lor \beta \leq_O \alpha$
\item \texttt{ord-zero-least}
  $\alpha \in \mathtt{ord} \;\Rightarrow\; 0 \leq_O \alpha$
\item \texttt{ord-lt-iff}
  $\alpha <_O \beta \;\Leftrightarrow\; \alpha \leq_O \beta \land \alpha \neq \beta$
\item \texttt{ord-succ-in}
  $\alpha \in \mathtt{ord} \;\Rightarrow\; \mathit{succ\_ord}(\alpha) \in \mathtt{ord}$
\item \texttt{ord-succ-above}
  $\alpha \in \mathtt{ord} \;\Rightarrow\; \alpha <_O \mathit{succ\_ord}(\alpha)$
\item \texttt{ord-succ-immediate}
  $\alpha,\beta \in \mathtt{ord},\; \alpha <_O \beta \;\Rightarrow\; \mathit{succ\_ord}(\alpha) \leq_O \beta$
\item \texttt{ord-succ-injective}
  $\mathit{succ\_ord}(\alpha) = \mathit{succ\_ord}(\beta) \;\Rightarrow\; \alpha = \beta$
\item \begin{sloppypar}\texttt{limit-ord-iff}
  $\mathit{limit\text{-}ord}(\lambda) \;\Leftrightarrow\; \lambda \in \mathtt{ord} \land \lambda \neq 0 \land \lnot\exists\alpha.\,\lambda = \mathit{succ\_ord}(\alpha)$\end{sloppypar}
\item \texttt{ord-segment-is-set}
  $\alpha \in \mathtt{ord} \;\Rightarrow\; \mathit{ord\text{-}segment}(\alpha) \in \mathtt{set}$
\item \texttt{ord-segment-membership}
  $x \in \mathit{ord\text{-}segment}(\alpha) \;\Leftrightarrow\; x <_O \alpha$
\item \texttt{ord-segment-zero}
  $\mathit{ord\text{-}segment}(0) = \emptyset$
\item \texttt{ord-segment-succ}
  \[x \in \mathit{ord\text{-}segment}(\mathit{succ\_ord}(\alpha))
  \;\Leftrightarrow\; x \in \mathit{ord\text{-}segment}(\alpha) \lor x = \alpha\]
\item \texttt{sup-ord-in}
  $A \in \mathtt{set},\; A \subseteq \mathtt{ord} \;\Rightarrow\; \mathit{sup\text{-}ord}(A) \in \mathtt{ord}$
\item \texttt{sup-ord-upper}
  $\mathit{sup\text{-}ord}(A)$ is an upper bound of $A$
\item \texttt{sup-ord-least}
  $\mathit{sup\text{-}ord}(A)$ is the least upper bound
\item \texttt{sup-ord-empty}
  $\mathit{sup\text{-}ord}(\emptyset) = 0$
\item \texttt{transfinite-induction}
  see \S\ref{sec:ord-induction}
\end{itemize}

Here $\leq_O$ and $<_O$ abbreviate \texttt{ord-le} and \texttt{ord-lt}.

\subsection{Relationship to number systems}

\texttt{nn} is a subset of \texttt{ord} (axiom \texttt{nn-subset-ord}):
\[
  \forall n.\; n \in \mathtt{nn} \;\Rightarrow\; n \in \mathtt{ord}
\]
Consequently $0 \in \mathtt{ord}$ (from \texttt{nn-zero-in} and
\texttt{nn-subset-ord}).  The ordinal successor agrees with the numeric
successor on $\mathtt{nn}$ (axiom \texttt{ord-succ-nn}), and
\texttt{ord-le} restricted to $\mathtt{nn}$ coincides with the numeric
\texttt{<=} (axiom \texttt{ord-le-nn-compat}).  The number systems
\texttt{zz}, \texttt{qq}, \texttt{rr}, \texttt{cc} are not subsets of
\texttt{ord}.

\subsection{Transfinite induction}
\label{sec:ord-induction}

\paragraph{The induction axiom.}

\begin{sloppypar}The axiom \texttt{transfinite-induction} is the \emph{complete} (strong)
form: a class $C$ that contains every ordinal whose strict predecessors are
all in $C$ must contain every ordinal.\end{sloppypar}
\[
  \forall C.\;
  \Bigl(\forall\alpha \in \mathtt{ord}.\;
        (\forall\beta <_O \alpha.\; \beta \in C) \Rightarrow \alpha \in C\Bigr)
  \;\Rightarrow\;
  \forall\alpha \in \mathtt{ord}.\; \alpha \in C
\]

\paragraph{The \texttt{(tfi)} proof command.}

\texttt{(tfi)} applies transfinite induction to the current focus goal.
The goal must have the form
\[
  \forall \alpha.\; \alpha \in \mathtt{ord} \;\Rightarrow\; P(\alpha)
\]
and \texttt{(tfi)} replaces it with the single inductive-step subgoal
\[
  \forall \alpha.\;
    \bigl(\alpha \in \mathtt{ord} \;\land\;
          \forall\beta <_O \alpha.\; P(\beta)\bigr)
    \;\Rightarrow\; P(\alpha)
\]

\paragraph{The \texttt{(tfi3)} proof command.}

\texttt{(tfi3)} splits the same goal into three subgoals matching the
standard Peano-style cases:
\begin{enumerate}
\item \textbf{Base:} $P(0)$.
\item \textbf{Successor:} $\forall\alpha \in \mathtt{ord}.\; P(\alpha) \Rightarrow P(\mathit{succ\_ord}(\alpha))$.
\item \textbf{Limit:} $\forall\lambda.\;
  \mathit{limit\text{-}ord}(\lambda) \land
  (\forall\beta <_O \lambda.\; P(\beta)) \;\Rightarrow\; P(\lambda)$.
\end{enumerate}

\subsection{Transfinite recursion}
\label{sec:ord-recursion}

The theoretical basis for transfinite recursive definitions is the
following pair of propositions.  Both are proved by the standard
least-counterexample argument using well-foundedness of $\leq_O$.

\begin{sloppypar}
\begin{prop}[Transfinite recursion, sets]\label{prop:tfrec-set}
Let $U = \mathit{ord\text{-}segment}(\alpha)$ for some $\alpha \in \mathtt{ord}$,
let $A \in \mathtt{set}$, and let
\[
  \mathcal{F}(U, A) \;=\; \bigcup_{x \in U}\,\mathtt{fun}(\mathit{ord\text{-}segment}(x),\, A)
\]
be the set of all functions defined on an initial segment within $U$ with
values in $A$.  For any $\mathcal{E} \in \mathtt{fun}(\mathcal{F}(U,A),\, A)$ there
is a unique $f \in \mathtt{fun}(U, A)$ such that
\[
  f(x) \;=\; \mathcal{E}\bigl(f \!\restriction\! \mathit{ord\text{-}segment}(x)\bigr)
  \quad \text{for every } x \in U.
\]
\end{prop}
\end{sloppypar}

\begin{prop}[Transfinite recursion, class version]\label{prop:tfrec-class}
Let $A$ be a class, and let $\mathcal{E}$ be a functoid from the class of
all functions defined on initial ordinal segments with values in $A$ to $A$
(a proper-class domain, so $\mathcal{E}$ cannot be an element of
$\mathtt{fun}$).
Then there is a unique functoid $f$ on $\mathtt{ord}$ with values in $A$ such that
\[
  f(\xi) \;=\; \mathcal{E}\bigl(f \!\restriction\! \mathit{ord\text{-}segment}(\xi)\bigr)
  \quad \text{for every } \xi \in \mathtt{ord}.
\]
\end{prop}

\begin{remark}
Both $\mathcal{E}$ and $f$ in Proposition~\ref{prop:tfrec-class} are functoids, not
elements of $\mathtt{fun}(A,B)$: their domains are proper classes ($\mathtt{ord}$
for $f$; the class of all initial-segment functions for $\mathcal{E}$), and
$\mathtt{fun}$ requires set-valued domains.  Each is nonetheless a well-defined
class of ordered pairs with the function property.
\end{remark}

The Scheme helper \texttt{def-by-ord-recursion} implements this pattern;
see Section~\ref{sec:recursive-defs}.

%% =========================================================
\section{The base library}
\label{sec:axioms}

The base library \texttt{VNB-LIBRARY} is loaded automatically and
provides the following named axioms (all usable via
\texttt{cmd-theorem-assumption}):

\begin{itemize}
\item \texttt{membership-implies-sethood}
  $\forall a\,\forall b.\;(a \in b) \Rightarrow a \in \mathtt{set}$
\item \texttt{extensionality}
  $\forall a\,\forall b.\;(a \in \mathtt{set} \land b \in \mathtt{set}) \Rightarrow
  (a = b \Leftrightarrow \forall x.\,(x \in a \Leftrightarrow x \in b))$
\item \texttt{empty-set-is-set}
  $\texttt{empty-set} \in \mathtt{set}$
\item \texttt{empty-set-has-no-members}
  $\forall x.\;\lnot(x \in \texttt{empty-set})$
\item \texttt{pairing}
  $\forall a\,\forall b.\;(a \in \mathtt{set} \land b \in \mathtt{set}) \Rightarrow \{a,b\} \in \mathtt{set}$
\item \texttt{pairing-membership}
  $\forall a\,\forall b.\;(a \in \mathtt{set} \land b \in \mathtt{set}) \Rightarrow
  \forall x.\;(x \in \{a,b\}) \Leftrightarrow (x=a \lor x=b)$
\item \texttt{power-set}
  $\forall a.\;a \in \mathtt{set} \Rightarrow \texttt{power}(a) \in \mathtt{set}$
\item \texttt{power-set-membership}
  $\forall a\,\forall x.\;(x \in \texttt{power}(a)) \Leftrightarrow
  (x \in \mathtt{set} \land \forall z.\,(z \in x \Rightarrow z \in a))$
\end{itemize}

A theorem is added by proving it.  State it with \texttt{sp}, prove it,
and install it with \texttt{qed}, which also installs its macete:
\begin{VNBcode}
(sp "forall([x in nn], x in nn)")
(di) (ass)
(qed 'my-thm)
\end{VNBcode}
\texttt{qed} refuses unless the proof is complete, and it is the only
command that records a result as proven.  A statement that is to be
assumed for the time being is added with \texttt{support} and given a
stated justification with \texttt{warrant!}
(Chapter~\ref{ch:classification}).  It is then recorded as asserted, and
every theorem that depends on it says so when it is installed.


%% Definitions moved to ch-defs.tex


%% =========================================================
\section{Pending axioms}
\label{sec:pending-axioms}

The axioms in this section are needed for full set-theoretic reasoning
but are not yet installed.  They fall into two groups.

\subsection*{Installed in \texttt{axioms.scm}}

The following were previously pending but are now installed.
They are straightforward first-order statements requiring no new
expression primitives.

\begin{itemize}

\item \texttt{equality-symmetry}
  $a = b \Rightarrow b = a$.

\item \texttt{equality-transitivity}
  $a = b \land b = c \Rightarrow a = c$.

\item \texttt{fun-apply-type}
  $(f \in \mathtt{fun}(A,B)) \land (x \in A) \Rightarrow f(x) \in B$.
  Derivable from \texttt{fun-codomain-iff} but named for direct use.

\item \texttt{list-sethood}
  $A \in \mathtt{set} \land L \in \mathtt{tuples}(A) \Rightarrow L \in \mathtt{set}$.
  Derivable from \texttt{subset-set} (elements of any class are sets)
  but named for direct use in structure instance proofs.

\item \texttt{eq-subst-membership}
  $a = b \land b \in S \Rightarrow a \in S$.
  A first-order instance of \texttt{equality-substitution} restricted to
  membership; derivable from \texttt{equality-transitivity} and
  extensionality, but named for direct use in type proofs.

\item \texttt{quasi-eq-reflexivity}\ $(a \mathbin{==} a)$.

\item \texttt{quasi-eq-symmetry}\ $(a \mathbin{==} b) \Rightarrow (b \mathbin{==} a)$.

\item \texttt{quasi-eq-transitivity}\
  $(a \mathbin{==} b) \land (b \mathbin{==} c) \Rightarrow (a \mathbin{==} c)$.

\item \texttt{quasi-eq-def}\
  $(a = a) \Rightarrow ((a \mathbin{==} b) \Leftrightarrow (a = b))$;
  when $a$ is defined, quasi-equality coincides with strict equality.

\end{itemize}

The following were already installed (noted here for completeness):
\texttt{fun-domain-apply-def}, \texttt{fun-domain-extensionality},
\texttt{fun-codomain-iff} and \texttt{choice-axiom} in \texttt{library.scm};
\texttt{nn-subset-ord} and all \texttt{ord-*} axioms in
\texttt{structure-library/ordinals.scm}, not in \texttt{library.scm}
(see Section~\ref{sec:ordinals}).  The ordinal axioms carry
\texttt{primitive} provenance and so contribute nothing to any proof's bill.

\subsection*{List recursion and list induction}
\label{sec:list-recursion}

Both were listed as pending in earlier editions of this manual, and both are
now installed.  \texttt{structure-library/list-recursion.scm} adds the
constructor \texttt{cons}$(x, L)$ --- the name is \texttt{cons}, not
\texttt{prepend} --- with four defining equations, stamped
\texttt{definitional} because they constrain a new head and so add no
consequence about the old vocabulary:
\[
  \texttt{length}(\texttt{cons}(x,L)) = \texttt{succ}(\texttt{length}(L)),
  \qquad
  \texttt{nth}(1, \texttt{cons}(x,L)) = x,
\]
with \texttt{nth-cons-succ} shifting the remaining indices and
\texttt{makeset-cons} relating \texttt{make-set} to a union with a singleton.
Alongside them are two \emph{generation} axioms, which are claims about
\texttt{tuples} rather than about \texttt{cons} and are therefore asserted and
warranted: \texttt{tuple-length-zero} (a tuple of length $0$ is the empty
tuple) and \texttt{tuple-cons-decompose} (a tuple of length $\texttt{succ}(n)$
is $x$ prepended to a tuple of length $n$).

The induction principle is then a \emph{theorem}, not a schema-level primitive
inference: \texttt{theorem-library/tuples-induction.scm} proves
\[
  \forall A, C.\;
  [\,]\in C \;\wedge\;
  (\forall x, M.\; x \in A \wedge M \in \mathtt{tuples}(A) \wedge M \in C
     \Rightarrow \texttt{cons}(x,M) \in C)
  \;\Rightarrow\;
  \forall L \in \mathtt{tuples}(A).\; L \in C
\]
by induction on the \emph{length}, using \texttt{nn-induction} and the two
generation axioms.  It is stated in class form, quantifying over a class $C$
with membership in it as the conclusion, exactly as \texttt{nn-induction} and
\texttt{finite-set-induction} are; VNB quantifies over classes, so one theorem
does the work a schema would do in a first-order presentation.  Its bill reads
\texttt{modulo \{tuple-length-zero, tuple-cons-decompose\}}, which is the honest
report: the induction is derived, and what it rests on is the statement that
tuples are finite.

\subsection*{Still pending}

The remaining items require schema-level infrastructure --- a primitive
inference rather than a first-order axiom.  (An earlier edition also listed
\texttt{separation}, \texttt{union-set}, \texttt{comp} and \texttt{prepend}
as missing expression primitives.  They are not missing: \texttt{sep},
\texttt{comp} and \texttt{big-union} are registered term-forming heads
(\texttt{wff.scm}) with intro/elim rules in \texttt{primitive-inferences.scm},
and they are documented in \texttt{reference/KERNEL-RULES.md}.)

\begin{itemize}

\item \texttt{equality-substitution}
  $a = b \land P(a) \Rightarrow P(b)$ (for any formula $P$).
  Requires a substitution schema; cannot be stated as a single
  first-order sentence.

\item \texttt{separation}
  $A \in \mathtt{set} \Rightarrow \{x \in A \mid p\} \in \mathtt{set}$.
  Requires the \texttt{separation} expression primitive.

\item \texttt{union-set} and \texttt{union-set-membership}
  $A \in \mathtt{set} \Rightarrow \bigcup A \in \mathtt{set}$, and
  $x \in \bigcup A \Leftrightarrow \exists a \in A.\; x \in a$.
  Requires the \texttt{union-set} (big union) expression primitive,
  distinct from the binary \texttt{union} constructor.

\item \texttt{cartesian-n-ary-sethood}
  All $A_i \in \mathtt{set} \Rightarrow A_1 \times \cdots \times A_n \in \mathtt{set}$.
  The binary case is installed (\texttt{cartesian-set-iff}); the $n$-ary
  case reduces to binary by associativity but requires explicit axioms.

\item \texttt{comp-membership}
  $x \in \{y \mid p\} \Leftrightarrow p[y \mapsto x]$.
  Requires the \texttt{comp} (class comprehension) expression primitive.

\item \texttt{lambda-type}
  $(\forall x_i \in A_i.\;\mathit{body} \in B) \Rightarrow
   \lambda([x_1,\ldots], \mathit{body}) \in \mathtt{fun}(A_1{\times}\cdots{\times}A_n, B)$.
  Requires a schema-level primitive inference (\texttt{pi-lambda-type}).

\item \texttt{lambda-beta}
  $\lambda([x_1,\ldots], \mathit{body})(v_1,\ldots,v_n)
   = \mathit{body}[x_i \mapsto v_i]$.
  Requires a beta-reduction primitive inference.

\item \texttt{iota-def}
  $(\exists!\, x.\;p(x)) \Rightarrow p(\texttt{iota}(x,p))$.
  Requires a schema-level primitive inference.

\item \texttt{infinity}
  There exists an inductive set:
  $\exists A \in \mathtt{set}.\;\emptyset \in A \land
   \forall x \in A.\; x \cup \{x\} \in A$.

\end{itemize}

%% =========================================================
